\chapter{\MakeUppercase{Ameaças à Validade}}\label{sec:ameacas}

Este capítulo discute as ameaças à validade do estudo segundo a taxonomia de \citeonline{Wohlin2012}, que distingue validade de conclusão, de construção, interna e externa. Para cada ameaça são descritos a origem, a mitigação adotada e o que permanece como limitação. Várias das ameaças foram identificadas durante o desenvolvimento e motivaram correções no pipeline; essas correções são registradas aqui porque documentam o que o protocolo final evita e o que ele ainda não evita.

\section{Validade de Conclusão}\label{sec2:amc-conclusao}

A validade de conclusão diz respeito à capacidade de inferir, com a estatística adequada, que existe relação entre o tratamento e o resultado.

\paragraph{Tamanho e tipicidade da amostra de decisões.} A ameaça mais séria do estudo é o número de decisões avaliadas. O teste cego tem 20 datas em duas safras, ambas de preço recorde; o consolidado com a validação tem 38 datas em quatro safras. Com essa amostra, a diferença entre o AgroVisible e a venda imediata não é significativa ao nível de 5\,\% no teste isolado e fica na margem ($p = 0{,}066$) no consolidado. Apenas a diferença contra a venda parcelada atinge significância ($p = 0{,}018$) com 38 datas. A mitigação foi ampliar a amostra ao período de validação, reportar intervalos de confiança por \textit{bootstrap} ao lado dos p-valores, e apresentar a fração do ganho do oráculo capturada, que é menos sensível ao nível de preços que a diferença em reais. A limitação permanece: as conclusões sobre ganho financeiro devem ser lidas como evidência favorável, não como demonstração.

\paragraph{Estrutura da diferença pareada.} O teste de Wilcoxon penaliza a política do motor por um motivo estrutural: na maioria das datas ele vende imediatamente e a diferença é exatamente zero, de modo que o ganho médio vem de poucas decisões com retorno alto. A fração de datas com diferença positiva contra a venda imediata é de 34\,\%, o que um leitor apressado tomaria por evidência de política ruim. A mitigação foi reportar, além do p-valor, a média, o intervalo de confiança e a decomposição por data, e explicar no texto a origem da mediana nula.

\paragraph{Variabilidade de treino da rede.} Diferenças de RMSE entre modelos podem ser ruído de inicialização. A mitigação foi treinar o modelo completo com cinco sementes e cada configuração da ablação com três, e interpretar apenas as diferenças maiores que o desvio entre sementes. Com esse critério, a vantagem da curva de futuros em 63 dias é robusta e a de 21 dias não é, e apenas dois grupos da ablação têm efeito interpretável no horizonte de 126 dias. A limitação é que três sementes por configuração ainda são poucas para efeitos pequenos, e a ablação deve ser lida como indicativa.

\paragraph{Múltiplas comparações.} O estudo compara sete estratégias, seis configurações de ablação e uma grade de sensibilidade, e não aplica correção para múltiplas comparações. As afirmações de significância referem-se a duas comparações pré-especificadas na RQ2, contra a venda imediata e contra a parcelada; as demais são descritivas.

\section{Validade de Construção}\label{sec2:amc-construcao}

A validade de construção diz respeito à aderência entre os conceitos que se pretende medir e as medidas efetivamente usadas.

\paragraph{Preço do produtor.} O sistema usa o indicador CEPEA/ESALQ como o preço que o produtor de Minas Gerais recebe. O indicador é a média de negócios nas praças de Minas e São Paulo, e o preço de balcão de uma cooperativa específica difere dele por um diferencial local que não é modelado. A consequência é que a base regional no sentido estrito, entre indicador e cooperativa, não é capturada; o que o sistema modela é a base entre o físico de referência e o futuro. Para um produtor cujo diferencial local é estável, a política é a mesma; para um cujo diferencial varia, há um risco de base adicional não contabilizado. Não há série pública que permita mitigar essa ameaça.

\paragraph{Instrumento de \textit{hedge}.} O contrato ICF da B3, referência doméstica, não tem série histórica pública com profundidade suficiente, e o contrato \textit{Coffee C} da ICE convertido pela PTAX foi usado como \textit{proxy}. O \textit{hedge} modelado não inclui margem de garantia nem ajuste diário, e assume caixa disponível para chamadas de margem sem descontar seu custo financeiro. O custo operacional de 10 R\$/sc foi tomado de material didático \cite{AgroSafety2026} sem decomposição em taxa de operação, emolumentos e registro. Como o motor nunca escolheu o \textit{hedge} no \textit{backtest}, essas simplificações não afetam os resultados de receita reportados, mas afetariam uma avaliação em que o \textit{hedge} fosse competitivo. No \textit{backtest}, a curva de futuros foi aproximada pelo primeiro vencimento porque a fonte de dados descarta contratos vencidos; a curva real por vencimento só passou a ser acumulada em setembro de 2026, e é ela que, na recomendação corrente, torna o \textit{hedge} de um mês a ação recomendada para aversão a risco moderada.

\paragraph{Custos de carregamento.} A tarifa de armazenagem de 2,50 R\$/sc/mês e a quebra técnica de 0,2\,\%/mês foram estimativas do desenvolvimento, sem documento de cooperativa ou armazém geral que as sustente. A mitigação foi a análise de sensibilidade, que mostrou receita e política insensíveis à armazenagem entre 1,50 e 5,00 R\$/sc/mês. A quebra técnica não foi variada. O custo de oportunidade usa a meta Selic como taxa livre de risco, sem considerar que o produtor pode ter custo de capital superior, por dívida de custeio, ou inferior, por aplicação a taxas de poupança; a sensibilidade entre 8 e 15\,\% a.a. cobre a faixa plausível e não altera decisões.

\paragraph{Custo de produção.} O COE por safra vem da planilha da CONAB para Guaxupé, sistema semiadensado de alta tecnologia e agricultura familiar, e não do custo real de um produtor específico. Duas safras dessa praça não têm COT e custo total na planilha. O COE entra no motor como piso opcional e como informação de margem, e a restrição de piso não atuou em nenhuma data do teste; portanto, a imprecisão do COE não afeta a receita reportada, apenas a classificação de margem exibida ao produtor.

\paragraph{Métricas preditivas.} RMSE e MAE são calculados sobre log-retornos e o MAPE sobre o preço reconstruído, o que os torna comparáveis entre horizontes mas não diretamente comparáveis com trabalhos que reportam erro sobre o preço em nível. A acurácia direcional foi acrescentada porque é a métrica que melhor antecipa o valor econômico, mas ela ignora a magnitude e trata um acerto de 1\,\% como igual a um de 30\,\%.

\paragraph{Heurísticas de comparação.} A venda parcelada foi implementada como seis parcelas mensais iguais, a média móvel como cruzamento da média de 63 dias e a meta de margem como 130\,\% do COE. Produtores reais usam variantes dessas regras com parâmetros distintos, e uma variante mais bem calibrada poderia reduzir a vantagem do motor. As heurísticas não foram otimizadas sobre os dados de teste, o que seria injusto com elas em sentido oposto.

\section{Validade Interna}\label{sec2:amc-interna}

A validade interna diz respeito à ausência de fatores não controlados que possam explicar o resultado no lugar do tratamento. Em um estudo de séries temporais, o fator dominante é o vazamento de informação futura.

\paragraph{Partição temporal.} A partição é estritamente cronológica por ano-safra, e o padronizador das variáveis é ajustado apenas no treino. Duas correções foram necessárias durante o desenvolvimento. Na primeira versão, as janelas de entrada eram construídas apenas dentro de cada conjunto, o que descartava os 62 primeiros dias da validação e do teste e fazia a primeira previsão partir de um preço quatro meses posterior ao último conhecido; a correção permitiu que a janela de entrada cruzasse a fronteira do conjunto anterior, por ser histórico já conhecido. Na primeira tentativa dessa correção o alvo também cruzou a fronteira, o que rotularia os últimos 126 dias do treino com preços da validação; isso foi detectado e corrigido antes do retreino, descartando como origem de amostra as últimas 126 datas de cada conjunto.

\paragraph{Data de disponibilidade das fontes.} O relatório de posições de fundos e os preços mensais de insumos foram inicialmente alinhados pela data de referência, o que usava informação ainda não publicada. Ambos foram reancorados na data de divulgação. Os acumulados de chuva foram calculados no calendário civil antes do recorte para dias úteis. As demais séries são diárias e alinhadas pelo último valor conhecido. Não há garantia de que nenhum vazamento sutil permaneça, mas os canais conhecidos foram auditados.

\paragraph{Estatísticas de base.} A base sazonal esperada e o risco de base foram, em uma primeira versão, estimados até o fim dos dados disponíveis, o que incluía a validação e o teste. A correção restringiu a estimativa ao período de treino. Essa correção foi feita antes dos experimentos reportados.

\paragraph{Viés no período de validação.} O \textit{backtest} estendido às safras 2022/23 e 2023/24 usa a dispersão residual e a parada antecipada calibradas nesses mesmos dados. O viés é favorável à rede e de magnitude pequena, porque a dispersão residual afeta apenas a largura do cone e não a direção da previsão, mas existe. Os resultados dessas duas safras são reportados separadamente dos do teste cego e o viés é declarado na tabela.

\paragraph{Modelo salvo desatualizado.} Durante o desenvolvimento, o caderno de previsão carregava o modelo do disco quando ele existia e chegou a produzir métricas com um modelo anterior à mudança nas janelas. Os experimentos reportados foram executados após retreino do zero, com semente fixa, e as métricas reproduziram exatamente as do dossiê de referência.

\paragraph{Seleção do modelo final.} A configuração da rede (unidades, dropout, janela, taxa de aprendizado) foi fixada a priori, sem busca em validação, o que evita sobreajuste ao conjunto de validação por seleção de hiperparâmetros. A semente 42 do modelo final ficou dentro de um desvio-padrão da média das cinco sementes em todas as métricas, o que afasta a hipótese de que o resultado seja um sorteio favorável.

\paragraph{Sobreajuste.} O RMSE de treino é de cerca de 0,03 contra 0,11 a 0,19 na validação, e a ablação mostra que remover três grupos de variáveis melhora o erro de longo prazo. O modelo completo está sobreajustado. Isso não invalida os resultados, que são medidos fora da amostra, mas indica que o ganho reportado é provavelmente um piso do que a arquitetura pode entregar.

\section{Validade Externa}\label{sec2:amc-externa}

A validade externa diz respeito à generalização dos resultados para além das condições do estudo.

\paragraph{Regime de mercado.} As quatro safras avaliadas cobrem uma queda de 40\,\%, uma recuperação a partir de um fundo, uma alta de 84\,\% e um platô alto. Faltam períodos de preço lateral prolongado e de preço baixo persistente, como 2018 a 2019, em que a política pode se comportar de forma diferente e em que a restrição de piso sobre o COE passaria a atuar. A geada de 2021 está no treino, não na avaliação. Uma avaliação em janela rolante ao longo de toda a série, retreinando a rede a cada safra, é o desenho que responderia a essa ameaça, e não foi executada por custo computacional e de tempo.

\paragraph{Outras praças cafeeiras.} O modelo foi treinado sobre o indicador do arábica em Minas e São Paulo, com clima de três regiões mineiras e paulistas. A transferência para o Espírito Santo, onde predomina o robusta, exigiria trocar o alvo pelo indicador do robusta, o instrumento de \textit{hedge} pelo contrato de conilon da B3 e as regiões climáticas, e reestimar a base. A arquitetura permite isso sem alteração de código, mas nenhum resultado deste trabalho se transfere diretamente.

\paragraph{Outras culturas.} A estrutura do motor de decisão é agnóstica à cultura: qualquer commodity armazenável com mercado a termo se ajusta ao mesmo processo de decisão. A transferência para grãos como milho e soja, que têm sazonalidade, custos de armazenagem e liquidez diferentes, exigiria novas fontes, novas variáveis e nova calibração, e a decisão de manter o escopo em uma única cultura foi tomada justamente porque cada cultura adicional multiplica esse trabalho.

\paragraph{Escala do produtor.} O \textit{backtest} avalia uma saca hipotética e assume que a venda não afeta o preço. Para um produtor pequeno essa hipótese é razoável; para uma cooperativa que comercializa milhões de sacas, o volume influencia a base local e a política precisaria de um modelo de impacto de mercado. O \textit{hedge} também é avaliado por saca, quando o contrato mínimo da B3 é de 100 sacas, o que exclui produtores muito pequenos da ação de proteção sem uma cooperativa intermediária.

\paragraph{Liquidez e execução.} O modelo assume que o produtor consegue vender a qualquer momento ao preço do indicador, sem custo de transação nem prazo de liquidação. Na prática, a venda física envolve classificação do lote, negociação com o comprador e prazo de pagamento, e o preço obtido pode diferir do indicador do dia. Essa hipótese de liquidez perfeita é comum na literatura e é retomada nas limitações do \autoref{sec:conclusao}.

\paragraph{Dependência de fontes de dados.} O indicador CEPEA exige download manual porque o portal bloqueia acesso automatizado; os contratos vencidos da ICE são descartados pela fonte gratuita; o contrato de robusta de Londres foi deslistado durante o projeto. Uma ferramenta em produção dependeria de acesso a dados pagos ou de acordo institucional, e a reprodutibilidade do estudo depende dos arquivos brutos preservados no repositório.
